\documentclass[11pt]{article}

\usepackage[margin=1in]{geometry}
\usepackage[T1]{fontenc}
\usepackage[utf8]{inputenc}
\usepackage{booktabs}
\usepackage{hyperref}
\usepackage{amsmath}

\title{Proactive Intent World Model}
\author{Anonymous ACL ARR Submission}
\date{}

\begin{document}
\maketitle

\begin{abstract}
Proactive retail assistance requires a model to infer a customer's latent state from visual behavior, compare possible interventions, and decide whether to act before the customer makes an explicit request. We propose the Proactive Intent World Model (PIWM), a vision-language framework for proactive retail guidance. PIWM couples visual customer-state perception with action-conditioned transition supervision derived from sales and consumer-behavior expertise. The current study organizes supervision into a general retail corpus and a compact target front-camera corpus for low-resource domain specialization. The resulting data suite supports state inference, next-state prediction, and five-act target-domain action selection over \textsc{Greet}, \textsc{Elicit}, \textsc{Inform}, \textsc{Recommend}, and \textsc{Hold}.
\end{abstract}

\section{Introduction}

Many vision-language agents are reactive: a user issues an instruction and the agent responds. Retail assistance requires a more proactive setting. A useful assistant must infer whether a customer is browsing, comparing, hesitating, or ready for help, and then choose whether to greet, elicit a need, provide information, recommend, or hold back.

PIWM treats proactive assistance as a world-modeling problem. Given short visual observations and a candidate action, the model predicts how the customer's state may change. This supports a policy that can compare counterfactual actions rather than directly imitating a single response.

\section{Method Overview}

PIWM uses a two-stage supervision design. Stage 1 trains general retail understanding through user-intent inference and action-conditioned next-state prediction. Stage 2 adapts the model to a target smart-vending/front-camera domain through five-act action selection. The operational five-act set is \textsc{Greet}, \textsc{Elicit}, \textsc{Inform}, \textsc{Recommend}, and \textsc{Hold}; \textsc{Reassure} is retained for source and compatibility analysis but excluded from the current operational action-selection path.

\section{Data}
\label{sec:data}

PIWM requires supervision that is different from standard visual question answering or dialogue-response data. A training example must connect a short visual observation to a customer state, a set of permissible proactive actions, and the expected consequence of each action. We therefore construct a \emph{general-to-target} data suite with two complementary domains. The general corpus teaches broad retail guidance from third-person, salesperson-observable scenes; the target corpus tests whether that guidance can specialize to a device-front-camera smart-vending setting with limited in-domain supervision.

\paragraph{Design principle: low-resource target specialization, not target-scale pretraining.}
The target corpus is intentionally compact. Its role is not to replace the general corpus with hundreds of additional videos, but to create a controlled domain shift: the model first learns proactive retail guidance in the general corpus, then adapts to a single fixed camera mounted on an intelligent vending terminal. This mirrors the deployment setting of a proactive service agent: broad sales knowledge is learned once, while each store, device, or product line provides only a small amount of domain-specific data. Our main data claim is therefore \emph{sample-efficient target specialization} rather than target-domain scale.

\subsection{From expert knowledge to structured outputs}

The PIWM outputs are not free-form labels written independently for each video. They are generated from an auditable expert-knowledge layer. We start from sales and consumer-behavior sources, including AIDA, OpenStax marketing chapters on consumer decision processes and personal selling, SPIN Selling, psychological reactance theory, and shopping-motivation work on hedonic, utilitarian, recreational, and gift-shopping behavior. Copyrighted sources are used only as citation-level anchors and compact paraphrases; the original long-form texts are not copied into the training set.

The knowledge base is organized in three stages. First, source passages are distilled into compact principles, such as when to ask an open need-focused question, when a recommendation should remain low-pressure, or when a firm recommendation risks psychological reactance. Second, these principles are promoted into conditional rules. The current corpus contains 27 distilled principles and 78 source-linked conditional rules, including cue-to-state priors, persona-to-intent rules, state-and-AIDA candidate-action rules, and action-conditioned transition rules. Third, the rules are compiled into deterministic output fields for each data record.

Concretely, a generated record follows the pipeline:
\[
\text{scene cues} \rightarrow
\text{customer state} \rightarrow
\text{intent tier} \rightarrow
\text{candidate actions} \rightarrow
\text{action outcomes} \rightarrow
\text{training targets}.
\]
The customer state includes the AIDA stage and a belief--desire--intention summary. The candidate-action rules produce a finite action set in the current five-act operational policy schema. The transition rules then produce, for each candidate action, a next state, reward, risk level, benefit level, and optional failure mode. The best action is the candidate with the strongest rule-supported outcome under a deterministic tie-break. Surface realizations are then filled from action templates: in the general corpus they correspond to human-salesperson guidance, while in the target corpus they correspond to terminal behavior such as screen text, voice style, light cue, or waiting state.

This construction matters for reviewability. If a label is wrong, the error can be traced back to a rule and its source link rather than to an opaque model output. It also gives counterfactual coverage: for the same visual state, the dataset can supervise not only the chosen action but also what would happen under inferior alternatives. This is the reason PIWM can train deliberation rows even when only one action is ultimately selected.

\subsection{Data sources and splits}

Table~\ref{tab:data-composition} summarizes the current data surface used by the EMNLP-style experiments. The general training split, \textsc{PIWM-Stage1-UserIntent-v1}, contains 543 synthetic video-backed parent states exported into 2,544 supervised examples: 543 user-intent rows and 2,001 action-conditioned next-state rows. The target source corpus contains 118 device-front-camera records from the smart-vending domain. For the main five-act experiment, we define the operational acts as \textsc{Greet}, \textsc{Elicit}, \textsc{Inform}, \textsc{Recommend}, and \textsc{Hold}. We exclude 17 records whose best act is \textsc{Reassure} and filter \textsc{Reassure} from candidate lists without creating empty candidate sets, yielding 101 clean five-act records. We use 71 records for target adaptation and reserve a balanced 30-record held-out target set with six records per operational act. This held-out target set is video-backed and project-lead QA-reviewed under the revised split.

\begin{table}[t]
\centering
\small
\begin{tabular}{llrrrrl}
\toprule
Split & Domain / view & Records & Frames & Examples & QA status & Role \\
\midrule
General train & Retail, third-person & 543 & -- & 2,544 & synthetic train & Stage-1 user state + next-state SFT \\
Target train & Smart-vending, frontcam & 71 & 213 & 71 & synthetic train & Stage-2 5-act adaptation \\
Target eval & Smart-vending, frontcam & 30 & 90 & 90 & project-lead QA-reviewed pass & balanced 5-act in-domain test \\
General eval & Retail, third-person & 36 & -- & 162 & QA-reviewed subset & forgetting check \\
\bottomrule
\end{tabular}
\caption{Current PIWM data composition. Target examples use three sampled frames per record. The main target experiment uses a balanced five-act split over \textsc{Greet}, \textsc{Elicit}, \textsc{Inform}, \textsc{Recommend}, and \textsc{Hold}; \textsc{Reassure} is excluded from the main macro-F1 and kept for source/compatibility analysis.}
\label{tab:data-composition}
\end{table}

\paragraph{General retail corpus.}
The general corpus covers ordinary in-store guidance from a salesperson-observable perspective. Each parent state is represented as a short visual window and a structured customer-state record. The supervision follows the PIWM decision loop: infer the current state, evaluate candidate interventions, and predict the next state and reward under each candidate action. This split preserves human-salesperson guidance logic and is used for Stage-1 training.

\paragraph{Target front-camera corpus.}
The target corpus comes from a fixed camera located at the intelligent vending terminal. This changes both the visual domain and the actor profile: the agent must infer customer intent from a device-front-camera view, and the response is a terminal behavior rather than a human salesperson's body movement. Each target record is sampled into three frames and mapped into the same v2.2 action schema as the general corpus. The main five-act train-test split is by record: 71 records are used for adaptation and 30 records are held out for target evaluation.

\subsection{Supervision format}

Each record is exported into one or more supervised examples, depending on which decision-loop component it trains. Table~\ref{tab:task-counts} reports the task decomposition. The user-intent task predicts the visible customer state, AIDA stage, intention label, and BDI summary from frames plus a plain scene description; it does not expose candidate actions, rewards, or best actions. The next-state task conditions on a candidate action and predicts the next AIDA stage and next BDI; \textsc{Hold}(\texttt{mode=silent}) rows serve as the no-intervention branch. The action-selection task asks the model to choose the best five-act intervention from candidate act--parameter realizations. Its prompt omits gold reward, risk, benefit, predicted next stage, and next state.

\begin{table}[t]
\centering
\small
\begin{tabular}{lrrrr}
\toprule
Split & User intent & Next-state prediction & Action selection & Total \\
\midrule
General train & 543 & 2,001 & 0 & 2,544 \\
Target train & 0 & 0 & 71 & 71 \\
Target eval & 30 & 30 & 30 & 90 \\
General eval & 36 & 126 & 0 & 162 \\
\bottomrule
\end{tabular}
\caption{Task-level decomposition of the revised two-stage training and evaluation files. Next-state rows are action-conditioned; the target evaluation next-state rows use the \textsc{Hold} no-intervention branch.}
\label{tab:task-counts}
\end{table}

\paragraph{Action schema.}
For the current target action-selection path, we report a five-act operational policy schema: \textsc{Greet}, \textsc{Elicit}, \textsc{Inform}, \textsc{Recommend}, and \textsc{Hold}. This is the action set used by the current Stage-2 target train/eval entrypoints. \textsc{Reassure} is excluded from the main macro-F1 and retained only as a source/compatibility boundary. Each act is parameterized rather than treated as a flat label; for example, \textsc{Recommend} distinguishes soft and firm pressure, \textsc{Greet} distinguishes open and close phases, and \textsc{Inform} distinguishes comparison, demo, attribute, and price content. This keeps the policy space small enough to audit while still preserving behaviorally meaningful variation.

\textsc{Reassure} is retained as a historical/source label and compatibility boundary. It is present in some target-front-camera source records, but it is excluded from the current five-act training, evaluation, inference, runtime-export, and macro-F1 path.

The target corpus is especially important for action coverage within this five-act space. It supplies device-front-camera examples for \textsc{Greet}, \textsc{Recommend}, and low-intervention \textsc{Hold} behavior, while the refreshed general v2 export represents \textsc{Recommend} with explicit soft/firm pressure parameters. No operational target act dominates the set: the largest raw class, \textsc{Inform}, accounts for 47 of 118 records.

\begin{table}[t]
\centering
\small
\begin{tabular}{lrr}
\toprule
Raw best act in target corpus & Records & Percent of 118 \\
\midrule
\textsc{Inform} & 47 & 39.8 \\
\textsc{Elicit} & 20 & 16.9 \\
\textsc{Greet} & 17 & 14.4 \\
\textsc{Reassure} (excluded from current 5-act) & 17 & 14.4 \\
\textsc{Recommend} & 11 & 9.3 \\
\textsc{Hold} & 6 & 5.1 \\
\bottomrule
\end{tabular}
\caption{Best-action distribution in the 118-record target front-camera corpus under the current five-act operational reporting schema. The 17 \textsc{Reassure} records are retained for compatibility/error analysis and are not counted in the current operational action-space claim.}
\label{tab:target-actions}
\end{table}

\subsection{Quality assurance}

We use QA as a data-split property rather than a blanket corpus label. For target evaluation, each held-out front-camera record is inspected through a contact sheet containing its three sampled frames. A record passes only if the customer is visible, gaze or head direction is interpretable, body posture is visible, terminal context is present, the three frames form a consistent temporal window, and the structured action label is supported by the visual evidence. The current 30-record held-out set was rebalanced into a balanced five-act split on May 19, 2026, and all 30 records passed project-lead QA after the new review templates were completed. The previous reviewed split is kept only as a historical audit artifact and is not used as the current target evaluation set.

This QA policy is deliberately conservative. The 30 target test records form the balanced in-domain evaluation set and are project-lead QA-reviewed under the revised split. The 71 target adaptation records remain synthetic training data: they are suitable for adaptation, but we do not describe the full 118-record target corpus as a fully human-reviewed benchmark. This separation prevents training-data scale from being confused with benchmark quality.

\subsection{Experimental use}

The data suite supports four comparisons. First, training on the general corpus and evaluating on the target held-out split measures zero-shot transfer from general retail guidance to the front-camera smart-vending domain. Second, continuing training on the 71 target five-act records measures low-resource target specialization. Third, training a joint model on the Stage-1 general rows plus the target five-act rows tests whether simple pooling matches staged specialization. Fourth, evaluating the target-specialized model back on the general QA set measures catastrophic forgetting. These comparisons are the reason the target split can be compact: its purpose is to expose a controlled domain shift with a reviewed in-domain test set after QA, not to serve as a large standalone pretraining corpus.


\section{Experiments}

The planned experiments compare general-only training, target-only training, staged general-to-target adaptation, and joint general-plus-target training. Evaluation focuses on macro-F1 and per-class F1 on a balanced target front-camera test set, plus a general-domain check for forgetting. The paper reports model size, training budget, package versions, split statistics, and action-label distributions.

\section{Limitations}

The current dataset remains compact, especially in the target front-camera domain. Some labels are derived from rule-based expert knowledge rather than large-scale naturally occurring interactions. Synthetic and staged videos may differ from real retail deployments, and proactive intervention can create privacy and over-intervention risks if deployed without additional review. The present study should therefore be interpreted as a controlled method and data-resource study rather than a claim of immediate deployment readiness.

\section{Ethics Statement}

The dataset design separates training data from reviewed evaluation data. Target evaluation records are manually reviewed for visual interpretability and label support. Any release of visual data should follow privacy review, consent handling, and de-identification procedures appropriate to the collection setting.

\end{document}
